\EMhold

\EMsection{معادلهٔ موجی دو‌بعدی}{معادلهٔ موجی دو‌بعدی}

\begin{EMunit}

مدلسازی یک غشاء مرتعش (به‌دست آوردن معادلهٔ حاکم بر \EMim{u(x,y,t)}).

\textbf{فرضیات فیزیکی:}

\begin{itemize}
\tightlist
\item
  جرم به صورت یکنواخت در طول غشاء توزیع شده (غشاء همگن)، کاملاً انعطاف‌پذیر است.
\item
  غشاء چنان محکم کشیده شده‌است که نیروی جاذبه در آن قابل چشم‌پوشی است.
\item
  میزان انحراف غشاء از حالت تعادل بسیار کوچک است.
\end{itemize}

\EMfigure{m06-membrane}{المان \EMim{\Delta x\times\Delta y} از غشاء مرتعش و نیروهای کششی
\EMim{T\Delta x} و \EMim{T\Delta y} وارد بر آن}

\end{EMunit}

\begin{EMunit}

با فرایند مدلسازی مشابه حالت یک‌بعدی، معادلهٔ ارتعاش غشاء دوبعدی به دست می‌آید:
\EMdisp{\frac{\partial^2u}{\partial t^2}\EMbr =c^2\left(\frac{\partial^2u}{\partial x^2}+\frac{\partial^2u}{\partial y^2}\right)\EMbr =c^2\nabla^2u}
برای حل یکتای این معادله نیاز به شرایط مرزی و شرایط اولیه داریم:

\begin{itemize}
\tightlist
\item
  شرایط مرزی: \EMim{u(x,y,t)=0} روی مرزهای غشاء
\item
  شرایط اولیه: \EMim{u(x,y,0)=f(x,y)} و \EMim{u_t(x,y,0)=g(x,y)}، که در آن \EMim{u_t\triangleq\dfrac{\partial u}{\partial t}}
\end{itemize}

\end{EMunit}

\EMhold

\EMsection{معادلهٔ موجی دو‌بعدی با شرایط مرزی مستطیل}{معادلهٔ موجی دو‌بعدی با شرایط مرزی مستطیل}

\begin{EMunit}

\EMfigure{m06-rectangle}{ناحیهٔ مستطیلی \EMim{0\le x\le a}، \EMim{0\le y\le b}}

\EMdisp{\frac{\partial^2u}{\partial t^2}\EMbr =c^2\left(\frac{\partial^2u}{\partial x^2}+\frac{\partial^2u}{\partial y^2}\right)}

\begin{itemize}
\tightlist
\item
  شرایط مرزی: \EMim{u(0,y,t)=u(a,y,t)=0} و \EMim{u(x,0,t)=u(x,b,t)=0}
\item
  شرایط اولیه: \EMim{u(x,y,0)=f(x,y)} و \EMim{u_t(x,y,0)=g(x,y)}
\end{itemize}

\end{EMunit}

\EMhold

\EMsubsection{حل معادله}{حل معادله}

\begin{EMunit}

در اینجا برای حل این معادله از روش \textbf{جداسازی متغیرها} استفاده می‌کنیم. فرض: \EMim{u(x,y,t)=X(x)\times Y(y)\times T(t)}.
\EMdisp{\frac{\partial^2u}{\partial x^2}\EMbr =X''YT,\qquad\EMbr \frac{\partial^2u}{\partial y^2}\EMbr =XY''T,\qquad\EMbr \frac{\partial^2u}{\partial t^2}\EMbr =XYT''\quad\EMbr \EMbr \Longrightarrow\quad\EMbr  XYT''\EMbr =c^2\big(X''YT+XY''T\big)}
\EMdisp{\frac{T''}{c^2T}\EMbr =\frac{X''}X+\frac{Y''}Y\EMbr =-\lambda^2\ \EMbr \Rightarrow\ \frac{Y''}Y\EMbr =-\lambda^2-\frac{X''}X\EMbr =-\rho^2\ \EMbr \Rightarrow\ \frac{X''}X\EMbr =-\lambda^2+\rho^2\EMbr =-\mu^2}
\EMdisp{\begin{cases}X''+\mu^2X=0\\ Y''+\rho^2Y=0\\ T''+c^2\lambda^2T=0\end{cases}\qquad\EMbr \lambda^2\EMbr =\mu^2+\rho^2}
\textbf{جواب در راستای \EMim{x}:}
\EMdisp{X(x)\EMbr =A_1\cos\mu x+A_2\sin\mu x,\quad\EMbr  X(0)\EMbr =0\EMbr \Rightarrow A_1\EMbr =0,\quad\EMbr  X(a)\EMbr =0\EMbr \Rightarrow A_2\sin\mu a\EMbr =0\EMbr \Rightarrow\mu a\EMbr =n\pi\EMbr \Rightarrow\mu_n\EMbr =\frac{n\pi}a}
\EMdisp{X_n(x)\EMbr =A_n\sin\mu_nx,\qquad\EMbr  \mu_n\EMbr =\frac{n\pi}a}
\textbf{جواب در راستای \EMim{y}:}
\EMdisp{Y(y)\EMbr =B_1\cos\rho y+B_2\sin\rho y,\quad\EMbr  Y(0)\EMbr =0\EMbr \Rightarrow B_1\EMbr =0,\quad\EMbr  Y(b)\EMbr =0\EMbr \Rightarrow B_2\sin\rho b\EMbr =0\EMbr \Rightarrow\rho b\EMbr =m\pi\EMbr \Rightarrow\rho_m\EMbr =\frac{m\pi}b}
\EMdisp{Y_m(y)\EMbr =B_m\sin\rho_my,\qquad\EMbr  \rho_m\EMbr =\frac{m\pi}b}
\textbf{جواب در راستای زمان:}
\EMdisp{\lambda_{nm}^2\EMbr =\mu_n^2+\rho_m^2\EMbr =\left(\frac{n\pi}a\right)^2+\left(\frac{m\pi}b\right)^2,\qquad\EMbr  T_{nm}(t)\EMbr =C_1\cos(c\lambda_{nm}t)+C_2\sin(c\lambda_{nm}t)}
\EMdisp{u_{nm}(x,y,t)\EMbr =\big[A_{nm}\cos(c\lambda_{nm}t)+B_{nm}\sin(c\lambda_{nm}t)\big]\sin\!\left(\frac{n\pi}ax\right)\sin\!\left(\frac{m\pi}by\right),\qquad\EMbr  n,m\EMbr =1,2,\dots}
\EMdisp{u(x,y,t)\EMbr =\sum_{n=1}^{\infty}\sum_{m=1}^{\infty}\big[A_{nm}\cos(c\lambda_{nm}t)+B_{nm}\sin(c\lambda_{nm}t)\big]\sin\!\left(\frac{n\pi}ax\right)\sin\!\left(\frac{m\pi}by\right)}

\textbf{اعمال شرط اولیه \EMim{u(x,y,0)=f(x,y)}:}
\EMdisp{f(x,y)\EMbr =\sum_{n=1}^{\infty}\sum_{m=1}^{\infty}A_{nm}\sin\!\left(\frac{n\pi}ax\right)\sin\!\left(\frac{m\pi}by\right)\EMbr =\sum_{n=1}^{\infty}K_n(y)\sin\frac{n\pi}ax,\qquad\EMbr  K_n(y)\triangleq\sum_{m=1}^{\infty}A_{nm}\sin\frac{m\pi}by}
\EMdisp{K_n(y)\EMbr =\frac2a\int_0^af(x,y)\sin\frac{n\pi}ax\,dx,\qquad\EMbr  A_{nm}\EMbr =\frac2b\int_0^bK_n(y)\sin\frac{m\pi}by\,dy}
\EMdisp{A_{nm}\EMbr =\frac4{ab}\int_0^b\left\{\int_0^af(x,y)\sin\frac{n\pi}ax\,dx\right\}\sin\frac{m\pi}by\,dy}

\end{EMunit}

\begin{EMexercise}{}

ثابت کنید:
\EMdisp{B_{nm}\EMbr =\frac4{abc\lambda_{nm}}\int_0^b\left\{\int_0^ag(x,y)\sin\frac{n\pi}ax\,dx\right\}\sin\frac{m\pi}by\,dy}

\end{EMexercise}

\begin{EMimportant}{جواب معادلهٔ موجی دو‌بعدی در مستطیل}

\EMdisp{u(x,y,t)\EMbr =\sum_{n=1}^{\infty}\sum_{m=1}^{\infty}\big[A_{nm}\cos(c\lambda_{nm}t)+B_{nm}\sin(c\lambda_{nm}t)\big]\sin\!\left(\frac{n\pi}ax\right)\sin\!\left(\frac{m\pi}by\right),\qquad\EMbr \lambda_{nm}^2\EMbr =\left(\frac{n\pi}a\right)^2+\left(\frac{m\pi}b\right)^2}
\EMdisp{A_{nm}\EMbr =\frac4{ab}\int_0^b\left\{\int_0^af(x,y)\sin\frac{n\pi}ax\,dx\right\}\sin\frac{m\pi}by\,dy,\qquad\EMbr  B_{nm}\EMbr =\frac4{abc\lambda_{nm}}\int_0^b\left\{\int_0^ag(x,y)\sin\frac{n\pi}ax\,dx\right\}\sin\frac{m\pi}by\,dy}

\end{EMimportant}

\begin{EMexample}{}

مطلوب است جواب معادلهٔ حرکت دو‌بعدی بالا با فرض \EMim{g(x,y)=0} و \EMim{f(x,y)=xy}.

\EMdisp{g(x,y)\EMbr =0\ \EMbr \Rightarrow\ B_{nm}\EMbr =0}
\EMdisp{A_{nm}\EMbr =\frac4{ab}\int_0^b\left\{\int_0^axy\sin\frac{n\pi}ax\,dx\right\}\sin\frac{m\pi}by\,dy\EMbr =\frac4{ab}\left\{\left[\int_0^by\sin\frac{m\pi}by\,dy\right]\times\left[\int_0^ax\sin\frac{n\pi}ax\,dx\right]\right\}}
\EMdisp{=\frac4{ab}\left[y\left(\frac{-b}{m\pi}\cos\frac{m\pi}by\right)-(1)\left(\frac{-b^2}{m^2\pi^2}\sin\frac{m\pi}by\right)\right]_0^b\times\left[x\left(\frac{-a}{n\pi}\cos\frac{n\pi}ax\right)-(1)\left(\frac{-a^2}{n^2\pi^2}\sin\frac{n\pi}ax\right)\right]_0^a}
\EMdisp{=\frac4{ab}\cdot\frac{-b^2}{m\pi}(-1)^m\cdot\frac{-a^2}{n\pi}(-1)^n\EMbr =\frac{4ab}{nm\pi^2}(-1)^{n+m}}
\EMdisp{u(x,y,t)\EMbr =\sum_{n=1}^{\infty}\sum_{m=1}^{\infty}\frac{4ab}{nm\pi^2}(-1)^{n+m}\cos(c\lambda_{nm}t)\sin\!\left(\frac{n\pi}ax\right)\sin\!\left(\frac{m\pi}by\right)}

\end{EMexample}

\EMhold

\EMsection{معادلهٔ گرمای دو‌بعدی با شرایط مرزی مستطیل}{معادلهٔ گرمای دو‌بعدی با شرایط مرزی مستطیل}

\begin{EMunit}

\EMdisp{\frac{\partial u}{\partial t}\EMbr =c^2\left(\frac{\partial^2u}{\partial x^2}+\frac{\partial^2u}{\partial y^2}\right)\EMbr =c^2\nabla^2u}

\begin{itemize}
\tightlist
\item
  شرط مرزی: \EMim{u(x,y,t)=0} روی مرزهای غشاء
\item
  شرط اولیه: \EMim{u(x,y,0)=f(x,y)}
\end{itemize}

\end{EMunit}

\EMsubsection{حل معادله}{حل معادله}

فرض: \EMim{u(x,y,t)=X(x)\times Y(y)\times T(t)}.
\EMdisp{\frac{\partial^2u}{\partial x^2}\EMbr =X''YT,\qquad\EMbr \frac{\partial^2u}{\partial y^2}\EMbr =XY''T,\qquad\EMbr \frac{\partial u}{\partial t}\EMbr =XYT'\quad\EMbr \EMbr \Longrightarrow\quad\EMbr  XYT'\EMbr =c^2\big(X''YT+XY''T\big)}
\EMdisp{\frac{T'}{c^2T}\EMbr =\frac{X''}X+\frac{Y''}Y\EMbr =-\lambda^2,\qquad\EMbr  \frac{Y''}Y\EMbr =-\lambda^2-\frac{X''}X\EMbr =-\rho^2,\qquad\EMbr \frac{X''}X\EMbr =-\lambda^2+\rho^2\EMbr =-\mu^2}
\EMdisp{\begin{cases}X''+\mu^2X=0\\ Y''+\rho^2Y=0\\ T'+c^2\lambda^2T=0\end{cases}\qquad\EMbr \lambda^2\EMbr =\mu^2+\rho^2}
مشابه حالت موجی (\EMim{X_n(x)=A_n\sin\mu_nx} با \EMim{\mu_n=\frac{n\pi}a} و \EMim{Y_m(y)=B_m\sin\rho_my} با \EMim{\rho_m=\frac{m\pi}b}) و با
\EMdisp{T_{nm}(t)\EMbr =C_1e^{-c^2\lambda_{nm}^2t},\qquad\EMbr \lambda_{nm}^2\EMbr =\mu_n^2+\rho_m^2\EMbr =\left(\frac{n\pi}a\right)^2+\left(\frac{m\pi}b\right)^2}
\EMdisp{u_{nm}(x,y,t)\EMbr =A_{nm}\sin\!\left(\frac{n\pi}ax\right)\sin\!\left(\frac{m\pi}by\right)e^{-c^2\lambda_{nm}^2t},\qquad\EMbr  n\EMbr =1,2,\dots,\ m\EMbr =1,2,\dots}
\EMdisp{u(x,y,t)\EMbr =\sum_{n=1}^{\infty}\sum_{m=1}^{\infty}A_{nm}\sin\!\left(\frac{n\pi}ax\right)\sin\!\left(\frac{m\pi}by\right)e^{-c^2\lambda_{nm}^2t}}
\textbf{اعمال شرط اولیه:} دقیقاً مانند حالت موجی
\EMdisp{A_{nm}\EMbr =\frac4{ab}\int_0^b\left\{\int_0^af(x,y)\sin\frac{n\pi}ax\,dx\right\}\sin\frac{m\pi}by\,dy}

\EMhold

\EMsection{بیان لاپلاسین در مختصات قطبی}{بیان لاپلاسین در مختصات قطبی}

\begin{EMunit}

هدف آن است که \EMim{\nabla^2u=u_{xx}+u_{yy}} در مختصات قطبی بر حسب \EMim{r,\theta} نوشته شود.

\EMfigure{m06-polar}{مختصات قطبی: \EMim{x=r\cos\theta}، \EMim{y=r\sin\theta}}

\end{EMunit}

\EMdisp{r\EMbr =\sqrt{x^2+y^2},\quad\EMbr  x\EMbr =r\cos\theta,\qquad\EMbr \theta\EMbr =\tan^{-1}\!\left(\frac yx\right),\quad\EMbr  y\EMbr =r\sin\theta}
\EMdisp{u_x\EMbr =u_rr_x+u_\theta\theta_x\ \EMbr \Rightarrow\ u_{xx}\EMbr =\big(u_rr_x+u_\theta\theta_x\big)_x\EMbr =u_{rr}r_x^2+u_{\theta\theta}\theta_x^2+2r_x\theta_xu_{r\theta}+r_{xx}u_r+\theta_{xx}u_\theta}
\EMdisp{u_y\EMbr =u_rr_y+u_\theta\theta_y\ \EMbr \Rightarrow\ u_{yy}\EMbr =u_{rr}r_y^2+u_{\theta\theta}\theta_y^2+2r_y\theta_yu_{r\theta}+r_{yy}u_r+\theta_{yy}u_\theta}
\EMdisp{r_x\EMbr =\frac xr,\quad\EMbr  r_{xx}\EMbr =\frac{y^2}{r^3},\quad\EMbr \theta_x\EMbr =\frac{-y}{r^2},\quad\EMbr \theta_{xx}\EMbr =\frac{2xy}{r^4},\qquad\EMbr  r_y\EMbr =\frac yr,\quad\EMbr  r_{yy}\EMbr =\frac{x^2}{r^3},\quad\EMbr \theta_y\EMbr =\frac x{r^2},\quad\EMbr \theta_{yy}\EMbr =\frac{-2xy}{r^4}}
\EMdisp{u_{xx}\EMbr =\frac{x^2}{r^2}u_{rr}+\frac{y^2}{r^4}u_{\theta\theta}+\frac{-2xy}{r^3}u_{r\theta}+\frac{y^2}{r^3}u_r+\frac{2xy}{r^4}u_\theta}
\EMdisp{u_{yy}\EMbr =\frac{y^2}{r^2}u_{rr}+\frac{x^2}{r^4}u_{\theta\theta}+\frac{2xy}{r^3}u_{r\theta}+\frac{x^2}{r^3}u_r+\frac{-2xy}{r^4}u_\theta}
با جمع این دو رابطه:

\begin{EMimportant}{لاپلاسین در مختصات قطبی}

\EMdisp{\nabla^2u\EMbr =u_{rr}+\frac1{r^2}u_{\theta\theta}+\frac1ru_r}

\end{EMimportant}

\EMhold

\EMsection{معادلهٔ موجی دو‌بعدی با شرایط مرزی دایره}{معادلهٔ موجی دو‌بعدی با شرایط مرزی دایره}

\begin{EMunit}

\EMfigure{m06-disc}{غشاء دایره‌ای به شعاع \EMim{R}}

\EMdisp{\frac{\partial^2u}{\partial t^2}\EMbr =c^2\nabla^2u,\qquad\EMbr  u(R,t)\EMbr =0\ \ \EMt{(شرط مرزی)},\qquad\EMbr  u(r,0)\EMbr =f(r),\ \ u_t(r,0)\EMbr =g(r)\ \ \EMt{(شرایط اولیه)}}

\end{EMunit}

\EMsubsection{حل معادله}{حل معادله}

با توجه به تقارن دایروی در این مسئله، بهتر است آن را در مختصات قطبی حل کنیم. بر روی غشاء دایروی، ارتعاش تابعی از \EMim{\theta} نیست، لذا داریم:
\EMdisp{\begin{cases}\dfrac{\partial^2u}{\partial t^2}=c^2\left(u_{rr}+\dfrac1ru_r\right)\\\noalign{\vskip 2mm} u(R,t)=0,\quad u(r,0)=f(r),\quad u_t(r,0)=g(r)\end{cases}}
حال با فرض تفکیک‌پذیری متغیرها \EMim{u(r,t)=P(r)\times T(t)} داریم:
\EMdisp{PT''\EMbr =c^2\left(P''T+\frac1rP'T\right)\ \EMbr \Rightarrow\ \frac{T''}{c^2T}\EMbr =\frac{P''}P+\frac1r\frac{P'}P\EMbr =-\lambda^2\ \EMbr \Rightarrow\ \begin{cases}P''+\dfrac1rP'+\lambda^2P=0\\\noalign{\vskip 1mm} T''+c^2\lambda^2T=0\end{cases}}

\EMhold

\EMsubsection{معادلهٔ دیفرانسیل بسل از مرتبهٔ \EMim{n}}{معادلهٔ دیفرانسیل بسل از مرتبهٔ n}

\begin{EMunit}

\EMdisp{y''+\frac1xy'+\left(1-\frac{n^2}{x^2}\right)y\EMbr =0\quad\EMbr \EMbr \Longrightarrow\quad\EMbr  y(x)\EMbr =AJ_n(x)+BY_n(x)}
که در آن \EMim{J_n} توابع بسل نوع اول و از مرتبهٔ \EMim{n} و \EMim{Y_n} توابع بسل نوع دوم و از مرتبهٔ \EMim{n} هستند.

\EMfigure{m06-bessel}{توابع بسل نوع اول \EMim{J_0,\dots,J_4} و نوع دوم \EMim{Y_0,\dots,Y_4}
برای \EMim{0<x\le10}}

\end{EMunit}

\begin{EMunit}

با تغییر متغیر \EMim{s=\lambda r}:
\EMdisp{P'\EMbr =\frac{dP}{dr}\EMbr =\frac{dP}{ds}\frac{ds}{dr}\EMbr =\lambda\frac{dP}{ds},\qquad\EMbr  P''\EMbr =\lambda^2\frac{d^2P}{ds^2},\qquad\EMbr \frac1r\lambda\EMbr =\frac{\lambda^2}s}
\EMdisp{\lambda^2\frac{d^2P}{ds^2}+\frac1r\lambda\frac{dP}{ds}+\lambda^2P\EMbr =0\ \EMbr \Rightarrow\ \frac{d^2P}{ds^2}+\frac1s\frac{dP}{ds}+P\EMbr =0}
این معادلهٔ دیفرانسیل بسل از مرتبهٔ صفر است (\EMim{n=0})، پس:
\EMdisp{P(r)\EMbr =AJ_0(\lambda r)+BY_0(\lambda r)}
\textbf{اعمال شرایط مرزی:}

\begin{itemize}
\tightlist
\item
  چون \EMim{P(0)\neq\infty} و \EMim{Y_0} در مبدأ بی‌کران است: \EMim{D=0} (ضریب \EMim{Y_0} صفر است) و \EMim{P(r)=CJ_0(\lambda r)}.
\item
  \EMim{P(R)=0} نتیجه می‌دهد \EMim{J_0(\lambda R)=0}، یعنی \EMim{\lambda R=\alpha_m} که در آن \EMim{\alpha_m} ریشه‌های \EMim{J_0} هستند:
  \EMdisp{\lambda\EMbr =\lambda_m\EMbr =\frac{\alpha_m}R,\quad\EMbr  m\EMbr =1,2,\dots\quad\EMbr \EMbr \Longrightarrow\quad\EMbr  P_m(r)\EMbr =C_mJ_0\!\left(\frac{\alpha_m}Rr\right)}
\end{itemize}

\EMfigure{m06-j0-zeros}{تابع \EMim{J_0(x)} و ریشه‌های آن \EMim{\alpha_1,\alpha_2,\alpha_3}}

\EMdisp{T''+c^2\lambda_m^2T\EMbr =0\ \EMbr \Rightarrow\ T_m(t)\EMbr =A_m\cos(c\lambda_mt)+B_m\sin(c\lambda_mt)}

\end{EMunit}

\begin{EMimportant}{جواب معادلهٔ موجی دو‌بعدی در دایره}

\EMdisp{u(r,t)\EMbr =\sum_{m=1}^{\infty}\big(A_m\cos(c\lambda_mt)+B_m\sin(c\lambda_mt)\big)J_0\!\left(\frac{\alpha_m}Rr\right),\qquad\EMbr \lambda_m\EMbr =\frac{\alpha_m}R}

\end{EMimportant}

\begin{EMremark}{تحقیق (تا ۱ نمرهٔ اضافه)}

گزارشی کامل از معادلهٔ دیفرانسیل بسل از مرتبهٔ \EMim{n} و توابع بسل نوع اول و دوم و خواص آنها و نیز حل کامل معادلهٔ دیفرانسیل موجی دوبعدی با شرایط مرزی دایره تهیه کنید (گزارش تایپ شده و در فرمت \lr{Word} باشد).

\end{EMremark}
